\subsubsection{Clustering Coefficient}
The package \verb+NetworkX+ has the function \verb+clustering+, with which the clustering of a graph can be calculated. 
The average shortest path length is calculated using the function \verb+nx.average_shortest_path_length+.
In our code we iterate over the different values of $p$ ($p$ = [0, 0.001, 0.003, 0.01, 0.03, 0.1, 0.3, 1]) as well as 1000 different seeds (1-1000). 
This way the results of our calculations are reproducible. 

The function \verb+networkx.clustering+ outputs a dictionary with the local clustering coefficients of each node. 
For each graph, we take the average of these local clustering coefficients and append them to a list for the current p. 
After iterating over all 1000 seeds, the average for the current p-value is taken and saved to another list.
Additionally, the standard deviation of the clustering coefficients is calculated using numpy and saved to yet another list. 
Then we continue with the next value of p. 

After iterating over all p values, the averaged values and their standard deviation are plotted using \verb+plt.errorbar+. 
All values are normalised by the value at $p=0$.
The whole code is timed, as we were afraid it would take too long to execute it.

To reproduce Figure \ref{fig:cluster_coeff} you need to run \verb+clustering_coeff.py+.
The required packages are \verb+NetworkX+, \verb+numpy+, \verb+matplotlib+ and \verb+time+.